\fichatitulo{01 · AVALIAÇÃO DE RAG}{Artigo de conferência · 2024}{ARES}
{\small SAAD-FALCON, Jon; KHATTAB, Omar; POTTS, Christopher; ZAHARIA, Matei.
\textit{ARES: An Automated Evaluation Framework for Retrieval-Augmented Generation Systems}.
In: NAACL-HLT 2024, v. 1: Long Papers. Mexico City: Association for Computational Linguistics, 2024. p. 338--354.
\href{https://doi.org/10.18653/v1/2024.naacl-long.20}{DOI: 10.18653/v1/2024.naacl-long.20}.\par}

\campo{Problema e objetivo}
Como avaliar componentes de RAG com menor esforço de anotação humana, preservando a capacidade de comparar configurações? O artigo propõe adaptar os avaliadores ao domínio e estimar a incerteza de seus julgamentos.

\campo{Principal contribuição · síntese interpretativa}
Combina juízes treinados com dados sintéticos e inferência apoiada por predições (PPI). Separa relevância do contexto, fidelidade da resposta e relevância da resposta; usa anotações humanas para corrigir estimativas e construir intervalos de confiança.

\campo{Método}
Treina juízes com dados sintéticos e avalia conjuntos de KILT, SuperGLUE e AIS; parte dos testes utiliza sistemas simulados com proporções conhecidas de acertos (seções 3--4).

\campo{Resultado principal}
Frente ao RAGAS 0.0.18, relata aumento médio de 0,065 e 0,132 no coeficiente de ordenação Kendall $\tau$, respectivamente para relevância do contexto e da resposta. São resultados da comparação experimental descrita, não ganhos de recuperação do nosso sistema (seção 5.1, p. 343, tabela 1).

\campo{Excertos-chave · seleção literal}
\begin{itemize}
\excerto{Dimensões}{context relevance, answer faithfulness, and answer relevance}{p. 338, Introdução; página 1 do PDF.}
\excerto{Calibração}{PPI improved the ranking prediction accuracy}{p. 343, seção 5.1; página 6 do PDF.}
\excerto{Participação humana}{ARES relies on a small set of annotations}{p. 346, seção 7; página 9 do PDF.}
\end{itemize}

\campo{Limitações declaradas e análise crítica}
\textbf{Autores:} necessidade de anotação especializada, recursos computacionais e avaliação principal em inglês (seção 7).
\textbf{Análise da ficha:} os testes com sistemas simulados e a mudança de idioma limitam a transferência direta para nosso caso.

\campo{Aproveitamento na dissertação · proposta}
\textbf{Destino:} seções 3 e 5; discussão de H3 na seção 8.
\textbf{Uso:} fundamentar a distinção entre dimensões avaliativas e a aferição do erro do juiz com referência humana. ARES não determina o tamanho da nossa amostra nem demonstra que H3 será corroborada.

\registro{Fonte consultada nesta preparação: PDF publicado, seções 1--7 (p. 338--346); apêndices não fichados. Excertos em inglês, sem tradução. Chave no projeto: \texttt{saadFalconEtAl2024ares}.}
